\section{The Liu--Rollenske parameter covering}

\begin{theorem}[The retained covering-parameter result]\label{ii:LR:covering}
\StatementRecord{ii:LR:covering}{proved}{The covering subclaim of the night-run compound statement P3K.36 was re-derived there by GPT-6 Astra (OpenAI, ChatGPT Pro mode, effort Pro), 16--18 November 2026. Also checked in this writing by \CurrentChecker: the exact lattice, affine quotient, Albanese map, intrinsic endpoint obstruction, rational square root and cap conjugations, and the two-step endpoint norm. The source and compactification constructions have their own records in Paper~IV.}{\ref{cdp:st:E1}; Paper~IV, \ref{night:P3F.41} and \ref{night:P3F.47}; the retained constructive maps of \ref{night:P3K.35} identified by \ref{night:P3K.74}; \cite[Lemma~2.7, Corollary~2.8, Proposition~2.12, Proposition~3.6 and Definition~3.7]{LiuRollenske2026}}
Retain the exact scalar $u$ in the Liu--Rollenske action $B_u=uB_0$,
the good fourth root at the $III^*$ source fibre, the specified
logarithmic vector $a=e_4/4$, and its fixed line character.
The resulting construction is a holomorphic family on the universal
logarithmic cover of a sufficiently small punctured $u$-disc.
The threefold reached after one positive $u$-turn is not
biholomorphic to the initial threefold. The actual family descends
to the degree-two parameter cover
\[
                         u=w^2.\tag{LR.1}
\]
This degree is minimal. The conclusion holds with every retained
constant root-line character.
\end{theorem}
\begin{proof}
\par

\medskip\noindent\textbf{1. The actual central torus and its scalar}

\par\nobreak

The source's equation (2.3) is the lift of translation by
\(-2P\), multiplied by \(u\). On the minimal good root the marked
section specializes to \((1+i)/2\) on
\(E=\mathbb C/(\mathbb Z+i\mathbb Z)\). Its line is the nonzero
order-two line with character
\[
                  m+in\longmapsto(-1)^{m+n}.          \tag{LR.4}
\]
These are the explicit statements and coordinate maps in
\cite[Lemma~2.7 and Corollary~2.8, pp.~6--7 of the attached-source edition]{LiuRollenske2026}. The same marked root and corrected line are constructed by the source and Cartier maps of Paper~IV, \ref{night:P3F.41} and \ref{night:P3F.47}.

Since \(2P\) specializes to zero, the restriction of \(B_0\) to
this line complement is multiplication by a constant \(c\ne0\).
The constant is independent of \(u\): every source, line, and root
comparison was fixed before multiplying the action by \(u\).
Thus its actual scalar multiplier is
\[
              \alpha=cu,\qquad
              \tau=\frac{\log(cu)}{2\pi i},
              \qquad\operatorname{Im}\tau>0.         \tag{LR.5}
\]
The line can be presented on the universal cover of \(E\) as
\((z,v)\sim(z+m+in,(-1)^{m+n}v)\). Writing
\(v=e^{2\pi i\eta}\), and then quotienting by \(B_u\), gives
exactly the central abelian surface
\[
 A_\tau=\mathbb C^2/\Lambda_\tau,
 \quad
 \Lambda_\tau=
 \left\langle
  (2,0),(1+i,0),(1,1/2),(0,\tau)
 \right\rangle_{\mathbb Z}.                           \tag{LR.6}
\]
The displayed lattice map is literal: its generators are the two
periods in the kernel of (LR.4), the remaining character period,
and the actual scalar quotient period. In particular
\[
 \delta=(0,1)=2(1,1/2)-(2,0),\qquad e_4=(0,\tau).
\]

The root deck acts on the source by \(z\mapsto-iz\). In this flat
line presentation an actual lift is
\((z,\eta)\mapsto(-iz,\eta+\kappa)\), with
\(4\kappa\in\mathbb Z\). Indeed (LR.4) is invariant under
\(-i\); the ratio of any two lifts is a nowhere-zero function on
the compact elliptic curve, hence a scalar; its fourth power is one.
The flat lift displayed in the proof of \cite[Proposition~3.9, p.~10 of the attached-source edition]{LiuRollenske2026} has \(\kappa=1/4\). Keeping \(\kappa\) records this actual lift and every retained constant finite line character. The logarithmic vector
is exactly \(a=(0,\tau/4)\), by \cite[Definition~3.7, p.~10 of the attached-source edition]{LiuRollenske2026}. Therefore the reduced multiple fibre is
\[
 B_{\tau,\kappa}=A_\tau/\langle R\rangle,
 \qquad R(z,\eta)=(-iz,\eta+\tau/4+\kappa).           \tag{LR.7}
\]
This action has order four and is free: a nontrivial power has a
nonintegral multiple of \(\operatorname{Im}\tau\) in its second
coordinate, whereas the imaginary second-coordinate periods of
\(\Lambda_\tau\) are integral multiples of
\(\operatorname{Im}\tau\). Its fourth power is the lattice period
\((0,\tau+4\kappa)\).

For reference, the linear deck in the ordered lattice basis of
(LR.6) is
\[
 D=\begin{pmatrix}
 1&1&0&0\\-2&-1&-1&0\\0&0&1&0\\0&0&0&1
 \end{pmatrix}.
\]
Its invariant lattice is exactly
\(\mathbb Z\delta\oplus\mathbb Z e_4\), agreeing with the
\cite[Proposition~3.6, pp.~9--10 of the attached-source edition]{LiuRollenske2026}.
A positive \(u\)-turn has the exact comparison
\[
 \tau\longmapsto\tau+1,
 \quad \Lambda_{\tau+1}=\Lambda_\tau,
 \quad a\longmapsto a+\delta/4.                      \tag{LR.8}
\]
The unchanged abelian surface in (LR.8) does not identify the affine
quotients: the torsion translation in (LR.7) has changed.

\par

\medskip\noindent\textbf{2. An intrinsic invariant of the reduced fibre}

\par\nobreak

The actual holomorphic projection
\[
 \pi_{\tau,\kappa}:B_{\tau,\kappa}\longrightarrow
 E_{\tau,\kappa}:=
 \frac{\mathbb C}{\frac12\mathbb Z+
                       (\frac\tau4+\kappa)\mathbb Z},
 \qquad [(z,\eta)]\longmapsto[\eta]                  \tag{LR.9}
\]
is well-defined by every generator in (LR.6) and (LR.7).
It has connected fibres. More explicitly, the image of the full
universal-cover deck group under its second-coordinate translation
is precisely the displayed rank-two lattice; after choosing a lift
of a target point, its zero-second-coordinate subgroup has elliptic
kernel generated by \((2,0),(1+i,0)\). This is the actual fibre of
(LR.9), so it is connected.

We also verify that (LR.9) is the Albanese map, without a choice of
marking. Holomorphic one-forms on \(A_\tau\) are constant linear
forms: the coefficients of their pullback to \(\mathbb C^2\) are
periodic entire functions and hence constant by compactness. The
\(R\)-invariant subspace is exactly \(\mathbb C\,d\eta\), since
\(R^*dz=-i\,dz\) and \(R^*d\eta=d\eta\). The periods of
\(d\eta\) on the full fundamental group of \(B_{\tau,\kappa}\)
are exactly \(\frac12\mathbb Z+(\frac\tau4+\kappa)\mathbb Z\):
this follows directly from the same affine deck generators. The
integration definition of the Albanese torus therefore gives
precisely (LR.9), up to a target translation. In particular its
elliptic-curve isomorphism class is intrinsic to \(B_{\tau,\kappa}\).

Normalize its first period to one by \(\eta\mapsto2\eta\). Its
Tate parameter is
\[
 Q=\exp\bigl(2\pi i(\tau/2+2\kappa)\bigr)
    =e^{4\pi i\kappa}\sqrt{cu}.                     \tag{LR.10}
\]
The factor \(e^{4\pi i\kappa}\) is a fixed sign. Under (LR.8)
this parameter changes to \(-Q\).

For completeness, the classical Weierstrass lattice invariant has
\(j(Q)=Q^{-1}+744+O(Q)\). The leading term alone suffices and can
be obtained directly from the convergent lattice sums: summing the
geometric series in the cotangent partial fraction gives
\[
 E_4(Q)=1+240\sum_{n\ge1}\sigma_3(n)Q^n,
 \quad E_6(Q)=1-504\sum_{n\ge1}\sigma_5(n)Q^n,
 \quad
 j(Q)=\frac{1728E_4(Q)^3}{E_4(Q)^3-E_6(Q)^2}.
\]
Their denominator is \(1728Q+O(Q^2)\), so \(Qj(Q)\) extends
holomorphically with value one at zero. Consequently
\[
 Q\bigl(j(Q)-j(-Q)\bigr)
\]
extends holomorphically with value two. It is nonzero on a
sufficiently small disc. Thus
\[
            B_{\tau,\kappa}\not\simeq
            B_{\tau+1,\kappa}                       \tag{LR.11}
\]
for every sufficiently small nonzero \(u\), even without preserving
any of their markings. A biholomorphism would induce an isomorphism
of their Albanese elliptic curves, whereas (LR.10) has different
Weierstrass invariant after the turn.

\par

\medskip\noindent\textbf{3. Every threefold endpoint must preserve this fibre}

\par\nobreak

We give the intrinsic argument independently of any canonical-bundle
formula. The actual source/cap Mayer--Vietoris calculation in Paper~IV, \ref{night:P3F.41} and \ref{night:P3F.47},
proves \(H^2(Y;\mathbb R)=0\) for these sphere members. Let
\(f:Y\to\mathbb P^1\) be their constructed connected torus fibration.
Every effective divisor on \(Y\) is vertical. Otherwise a horizontal
irreducible component meets a general smooth torus fibre \(F\) in
a nonzero effective divisor \(C\). Its real first Chern class is
zero, because it is the restriction of the corresponding class on
\(Y\). Every compact complex torus admits a translation-invariant
Kähler form \(\omega\), and the divisor identity gives
\(\int_F c_1(\mathcal O(C))\wedge\omega=\int_C\omega>0\), a
contradiction.

Let \(g\) be a meromorphic function on \(Y\). Its poles are
therefore vertical. On a general fibre it has no divisorial poles;
normality and Hartogs extension remove the remaining codimension-two
indeterminacy. It is a holomorphic function on a compact connected
torus and is constant. This gives a meromorphic function on the
regular part of the base. Near a bad base value \(p\), multiply
\(g\) by a sufficiently large integer power of \(f^*(t-p)\) to
cancel every polar order on the finitely many fibre components.
The product is holomorphic by normality. It is constant on compact
connected fibres and hence descends holomorphically: its descended
function is bounded near \(p\) by properness, and the removable
singularity theorem applies. Dividing back proves
\[
                   \mathbb C(Y)=f^*\mathbb C(t).      \tag{LR.12}
\]

A biholomorphism between two such threefolds induces an automorphism
of this meromorphic function field, hence a Möbius map of the base.
Indeed a rational function with a rational compositional inverse
has degree one. Thus it fits into an actual commutative square with
their maps \(f\). Multiplicities of fibres are preserved, because a
Möbius map has nonzero local derivative. The unique multiple fibre
has multiplicity four, so its reduced support must map biholomorphically
to the unique reduced multiple fibre in the target. This contradicts
(LR.11) after a single positive \(u\)-turn. This argument excludes
all possible holomorphic endpoint maps, including source/base
symmetries and source inversion; no assumed translation normal form
for an endpoint is needed.

\par


\medskip\noindent\textbf{4. The actual map for the double turn.}\par\nobreak
We display the retained construction, including its scalar norm.
Write $G=P$ for the marked source section and
$\mathcal M=\mathcal O(G-O)$ for the line used in the source.
On the generic elliptic fibre the rational function
\[
 f=\frac1{x-x(G)},\qquad
 \operatorname{div}(f)=O+O-G-(-G)
                       =t_G^*(G-O)-(G-O)
\]
gives a line isomorphism covering translation by $+G$.
Let $U$ be its meromorphic bundle lift, with the prescribed
vertical Cartier corrections retained. Since $B_0^{-1}$
covers translation by $+2G$, their exact ratio is a rational
base function
\[
              \rho(t)=B_0^{-1}/U^2.\tag{LR.13}
\]
Indeed the ratio is a nonvanishing holomorphic function on the
generic compact elliptic fibre, hence constant there; the
explicit divisor maps make that constant rational in $t$.
At the selected $I_1$ value, write $e$ for the vertical order
of $U$. The actual linearization has a simple zero there, so
\[
                  \operatorname{ord}\rho=-1-2e.\tag{LR.14}
\]
At an uncapped retained fibre the corresponding order is $-2e$.
The integral component-intersection kernel is generated by the
full reduced fibre; thus removing that integral whole-fibre
power makes $U$ invertible there.

On the good fourth root let $z$ be the source-root coordinate,
let $D$ be its deck, and let $L$ be the extending lift of $+G$.
The source and line correspondence of Paper~IV constructs this
lift on the actual corrected root line. The point $G(0)$ has
order two. The fibre weights of $\mathcal O(G-O)$ at $G(0)$
and $O(0)$ are the deck derivative and its inverse. Their
ratio is $-1$, since the derivative has order four. A retained
constant character multiplies both weights by the same number
and leaves the ratio unchanged. Hence
\[
                 DLD^{-1}=-L.\tag{LR.15}
\]
Initially the commutator may be $-c(z)$ with $c(0)=1$ and
$\prod_{j=0}^3c(i^jz)=1$. Put $g=\log c$ and
$\mathcal D h(z)=h(-iz)$. The explicit function
$h=-\tfrac14\sum_{j=0}^3j\mathcal D^jg$ satisfies
$(\mathcal D-I)h=-g$. Rescaling the target of $L$ by $e^h$
gives (LR.15). The unit $L^2/B_0^{-1}$ is then deck invariant;
its invariant analytic square root further normalizes
\[
                              L^2=B_0^{-1}.\tag{LR.16}
\]
This rescaling preserves (LR.15).

Write $U=z^eh(z)L$ on the root line, with $h$ a nonzero unit.
Because $U$ is defined over the original base, its deck
conjugation gives $e\equiv2\pmod4$ and $h(iz)=h(z)$.
In particular
\[
                    \operatorname{ord}_{t=0}\rho=-e/2
                         \equiv1\pmod2.\tag{LR.17}
\]
Let $h_0(t)$ be the actual positive one-gamma closing multiplier
of the degree-four parameter family $u=r^4$. It has order
$-1$ at this cap, order $+1$ at the selected $I_1$, and no
other orders except the retained ordinary framing orders,
which are divisible by four. Equations (LR.14) and (LR.17)
show that $h_0\rho$ has even order at every point of
$\mathbb P^1$. Factoring its zeros and poles and choosing the
specified cusp value therefore gives a rational function
$g_0$ with
\[
                             g_0^2=h_0\rho.\tag{LR.18}
\]
The actual ordinary-chart endpoint covering $r\mapsto-r$ is
\[
                    \widetilde D_r=-r^{-2}g_0(t)U.\tag{LR.19}
\]
The minus sign is the scalar norm correction and is part of
this map.

On the good-root cap the apparent $z^{-2}$ factor in
(LR.19) is cancelled by the actual ordinary-to-cap gluing
change under $r\mapsto-r$. The remaining map is an invariant
unit times $L$, so it extends invertibly through the cap.
Writing its action as $A_r=rA_0$, equation (LR.15) gives
\[
               LA_rL^{-1}=-A_r=A_{-r}.\tag{LR.20}
\]
It also commutes with $B_u$, since $(-r)^4=r^4$.
At the retained order-one datum the deck is the identity,
and the ordinary comparison extends directly. These are
maps of the actual quotients, with their finite characters.

The cusp comparison includes both source and line periods.
Put $v=\log r/(2\pi i)$ and write the actual bounded period
expressions as
\[
 H=s+U_0(q),\qquad
 \gamma=(\Gamma,C-s-4v),\qquad
 \tau(v)=\frac{H\delta+\gamma}{2}.
\]
The endpoint lift is $\tau(v)+\delta/2$.
Its two steps satisfy the literal identity
\[
 [\tau(v)+\delta/2]+[\tau(v+1/2)+\delta/2]
                       =H\delta+\gamma.\tag{LR.21}
\]
Every cusp Cox character is multiplied by its bounded
holomorphic unit; this extends through every retained cone
and face of the chosen periodic fan. On the ordinary chart
the identical norm is
\[
 \widetilde D_{-r}\widetilde D_r
    =r^{-4}g_0^2U^2
    =h_0 B_u^{-1}.\tag{LR.22}
\]
Thus the second iterate is exactly the prescribed one-gamma
gluing, including the nonconstant multiplier $h_0$.

First quotient the universal $v$-family by this one-gamma
map to obtain the proper $r$-family. Equation (LR.21) then
makes (LR.19) an order-two action on that family, covering
$r\mapsto-r$. It is free because $r\ne0$. Its quotient
is a proper holomorphic family on the punctured $w$-disc,
where $w=r^2$ and $u=w^2$. The cap and cusp extensions above
prove smoothness of its total space. Steps 2--3 exclude
every biholomorphic endpoint for one positive $u$-turn, so
the degree two is minimal.
\end{proof}

\begin{refuted}[The historical all-fan clause of P3K.36]\label{night:P3K.36}\label{stmt:P3K.36}
\StatementRecord{night:P3K.36}{refuted compound claim; covering subclaim retained as \ref{ii:LR:covering}}{The source itself records its refutation by P3K.74. Current scoped refutation review by \CurrentChecker: the complete historical block and P3K.74 counterexample were read; for the identical endpoint $u=w^2$, the specified $T_2$ datum has actual value $26$ in $\mathbb Z/28$ while the withdrawn polynomial gives zero. The actual covering proof is recorded separately at \ref{ii:LR:covering}. The underlying full fan proof retains its own current check record in Paper~IV.}{Paper~IV, \ref{night:P3K.73} and \ref{night:P3K.74}}
The earlier clause assigned the value
$3(S-1)(S-3)/8\pmod {28}$ to every fan of the corrected
$u=w^2$ family, where $S=p^2-pq+q^2$ was the norm of the
physical delta cocharacter. The actual fan $T_2$ has $S=1$,
so that expression gives zero. The retained degree-four
page has $(H,\chi)=(2,1)$ and its degree-two quotient has
$(J,\widehat H,\widehat\chi)=(2,1,1/2)$.
The actual fan maps in Paper~IV, \ref{night:P3K.73} and
\ref{night:P3K.74}, give
\[
 (\widehat\chi-\widehat H/2)-J\,2(2-1)/2
                           =-2=26\pmod {28}.
\]
This is the stated counterexample to the all-fan clause.
The original fixed reference fan retains its separately
proved zero value; the single $u$-turn has no endpoint
identification by \ref{ii:LR:covering}.
\end{refuted}
